\chapter{Introduction}

\section{Biological Context for Computer Scientists}
To contextualize the computational challenge of \gls{dtu}, it is necessary to establish the physical relationship between genes and their transcribed RNA\defn{RNA (ribonucleic acid)}{The broad class of nucleic acid molecules transcribed from DNA. It comprises messenger RNA (mRNA), which carries protein-coding instructions, and non-coding RNAs that perform structural or regulatory roles.} outputs. While all cells in an organism share an identical DNA genome, cellular diversity is driven by which specific segments of that genome are actively expressed. A gene serves as the baseline sequence template. When a gene is activated, its genetic information is copied into a temporary messenger molecule called a transcript (messenger RNA, or mRNA)\defn{messenger RNA (mRNA)}{The protein-coding subtype of RNA that carries a gene's sequence to the ribosome; it is the RNA species measured by scRNA-seq.}. Crucially, the mapping from a parent gene to a mature transcript is rarely one-to-one. A single gene is physically composed of coding blocks (exons) separated by non-coding boundaries (splice junctions). During RNA maturation, the cell can selectively include or exclude specific exons---a combinatorial process known as \gls{splicing}. This mechanism allows a single parent gene to generate multiple distinct structural variants, or \gls{isoform}. Because distinct isoforms frequently possess divergent or even antagonistic functional properties, a cell can fundamentally alter its biological behavior simply by shifting which specific isoform it produces---an alternative MET transcript (MET$\Delta$14), for instance, drives ligand-dependent invasive growth in cancer \cite{cerqua2022met}. This can occur even if the total, aggregated expression level of the parent gene remains strictly constant. Detecting this targeted internal shift is the core objective of DTU analysis.

\section{The Main Technical Problem}
Historically, \gls{rnaSeq} aggregated mRNA \glspl{read} across entire populations, in a protocol known as \gls{bulkRnaSeq}. Bulk profiling averages the RNA content of many cells at once, effectively erasing cellular heterogeneity\defn{Cellular heterogeneity}{The coexistence of distinct cell types and states within a single organism.}. \gls{scRnaSeq} resolved this by profiling thousands of individual cells simultaneously, capturing the dynamic expression landscapes unique to each cell. Yet modern scRNA-seq relies on a pervasive hardware compromise. To economically profile tens to hundreds of thousands of cells, widely adopted droplet protocols\defn{Droplet protocols}{A high-throughput single-cell method that encapsulates each cell in a tiny oil droplet with a barcode, enabling massive-scale profiling at the cost of capturing mostly transcript ends.} use \gls{threePrime}, capturing only the \gls{threePrimeEnd} (the extreme tail) of each mRNA molecule. Conversely, \gls{fullLength} maps the entire transcript but is restricted by labor-intensive logistics and high costs to merely hundreds or thousands of cells.

Because 3'-biased data is enough to accurately identify the parent gene, it has cemented Differential Expression (DE)\defn{Differential Expression (DE)}{The statistical comparison of a gene's total abundance between conditions; blind to which isoform is produced.}---the statistical comparison of total gene abundance---as the cornerstone of transcriptomic analysis. This emphasis is well justified: much of the functional specialization that distinguishes cell types and states arises at the level of gene regulation, as cells selectively activate or silence entire genes. DE captures precisely this layer, reliably isolating upregulated or downregulated disease markers\defn{Marker genes}{A gene or transcript whose expression specifically identifies a cell type, state, or condition.} and delivering the bulk of the signal in most transcriptomic studies. It does, however, operate with a major biological blind spot: it assumes total transcript abundance dictates function. If a cell switches from producing one isoform to a functionally distinct variant while keeping total gene expression constant, DE registers this critical regulatory event as ``no change''. DTU analysis addresses this by tracking shifts in relative isoform proportions, independent of overall gene expression.

Scaling DTU to massive droplet-based datasets remains a fundamental analytical bottleneck. Short 3' reads inherently lack coverage of internal splice junctions, the exact locations where alternative splicing occurs. This lack of structural resolution forces analysts to sum reads at the gene level, completely obscuring isoform dynamics and leaving a vital regulatory layer hidden within standard scRNA-seq data.

\section{The Computational Challenge}
Recovering \gls{dtu} from 3'-biased \gls{scRnaSeq} data presents two formidable computational hurdles:

\begin{itemize}
    \item \textbf{Extreme \gls{sparsity}:} Single-cell count matrices are overwhelmingly sparse, with zero entries routinely exceeding 70\% to 90\% of all values. This sparsity is a confounding mixture of true biological absence and \gls{dropout}---false absences in which a transcript is genuinely present but missed because capture or sequencing failed. Because standard statistical models assume continuous distributions, they cannot disentangle these distinct zero types, leading to fundamentally flawed estimates when applied to such sparse data.

    \item \textbf{Positional Bias:} High-throughput droplet assays capture only the terminal ends of transcripts, meaning the resulting data inherently lacks the internal read coverage required to disambiguate structural variants. When multiple complex \glspl{isoform} share identical 3' terminal exons, pseudoalignment tools\defn{Pseudoalignment}{A fast mapping strategy that assigns a read to candidate transcripts by matching short $k$-mer substrings instead of base-by-base alignment.} suffer from massive read-assignment ambiguity, as the reads carry no structural information to distinguish between the variants.
\end{itemize}

Overcoming these barriers requires a novel statistical framework---one capable of rescuing genuine regulatory signals from heavily \glspl{zeroInflated} while successfully navigating the ambiguity of incomplete transcript coverage.

\section{Current Approaches and Limitations}
Existing computational strategies attempt to recover transcript-level regulation through several distinct paradigms, yet each introduces severe methodological compromises:

\begin{itemize}
    \item \textbf{Full-Length Protocols:} Workflows built on full-length scRNA-seq (such as SMART-seq) capture complete transcript bodies, enabling robust isoform-switching analyses with tools like IsoformSwitchAnalyzeR\defn{IsoformSwitchAnalyzeR}{A bioinformatics package designed to detect isoform switching and predict functional consequences from full-length RNA-seq data, requiring complete transcript read coverage.}. However, their reliance on plate-based logistics and high sequencing costs restricts throughput to hundreds or a few thousand cells, preventing population-scale screening.

    \item \textbf{Bulk-Adapted Algorithms:} Established bulk RNA-seq tools (including rMATS, SUPPA2, DRIMSeq, and DEXSeq) rely on continuous isoform proportion modeling and deep junction-spanning read depth. When applied directly to sparse single-cell matrices, their statistical assumptions break down, producing degenerate estimates and high error rates.

    \item \textbf{Pseudo-Bulk Aggregation:} To restore the coverage required by bulk algorithms, researchers often perform \gls{pseudoBulk}, aggregating single cells into ``pseudo-bulk'' profiles. While computationally tractable, this pooling discards the cell-to-cell variance that single-cell sequencing was fundamentally designed to resolve.
    
    \item \textbf{Targeted 3'-Biased Tools:} Specialized frameworks designed for droplet data (such as Sierra and MANTA) focus primarily on a narrow class of structural events defined at transcript termini. They do not provide a generalized, genome-wide framework for detecting multi-isoform switching from raw zero-inflated count matrices.
\end{itemize}

Consequently, the field currently lacks a sparsity-robust statistical framework capable of resolving Differential Transcript Usage directly at single-cell resolution from standard droplet-based experiments.

\section{Thesis Objective}
The primary objective of this thesis is to determine whether sparsity-robust statistical modeling can overcome the droplet sequencing bottleneck to detect Differential Transcript Usage at true single-cell resolution.

This investigation builds upon COTAN (Co-expression Tables Analysis)\defn{COTAN (Co-expression Tables Analysis)}{A statistical framework designed for zero-heavy single-cell data that models non-zero and zero counts directly via $2 \times 2$ contingency tables rather than continuous distribution approximations.}, a mathematical framework originally developed to model zero counts in gene-level \gls{coexpression} via presence/absence contingency tables. Given COTAN's ability to extract biological signals from sparse gene matrices without artificial imputation\defn{Imputation}{Replacing zero counts with values predicted from neighbouring cells, at the risk of inventing spurious correlations.} or variance-stabilizing log-transformations, this work tests a central hypothesis: \textbf{can contingency-table independence testing be successfully extended from gene co-expression to transcript-level isoform mutual exclusivity?}

To address this challenge, this thesis delivers three core contributions:

\begin{enumerate}
    \item \textbf{Algorithmic Adaptation:} Formulating an isoform-aware detection framework that identifies mutually exclusive transcript pairs within individual genes, followed by bidirectional contrast filtering across cell clusters to isolate genuine biological switches.

    \item \textbf{A Reproducible Quantification Pipeline:} Engineering an automated, containerized Nextflow\defn{Nextflow}{A workflow language that runs each analysis step in its own container, making pipelines portable and reproducible.} workflow utilizing modified \texttt{alevin-fry}\defn{alevin-fry}{A fast single-cell quantification tool that assigns reads to transcripts by pseudo-alignment.} indexing (identity transcript-to-gene mapping) to derive transcript-level count matrices directly from raw 3'-biased reads without gene-level collapse.

    \item \textbf{Empirical Validation:} Benchmarking the adapted framework against both homogeneous human cell line mixtures and complex mouse cortical tissue, cross-referencing candidate switches against published isoform biology and comparing transcript-derived clusters against gene-level baselines.
\end{enumerate}

This thesis does not claim to circumvent the physical constraints of 3'-biased chemistry; isoforms with identical 3' terminal exons remain physically indistinguishable. Instead, this work establishes the practical boundary of computational recovery: demonstrating that wherever terminal sequence divergence exists, sparsity-robust contingency statistics can reliably recover differential transcript usage across massive droplet-based atlases without the prohibitive expense of full-length sequencing.

\section{Thesis Outline}
The remainder of this thesis is structured as follows.

Chapter 2 (Theory and Frameworks) establishes the computational and biological foundations for transcript-level single-cell analysis. It details the mechanics of pseudo-alignment and Expectation-Maximization for multi-mapping reads, examines the statistical hurdles of extreme matrix sparsity, and introduces the contingency-table mathematics underlying the COTAN framework.

Chapter 3 (Materials and Methods) documents the reproducible methodology and dataset selection criteria used throughout this work. It details the custom Nextflow pipeline engineered to construct transcript-level count matrices via \textit{alevin-fry}, and explicitly breaks down the step-by-step DTU detection algorithm adapted from COTAN.

Chapter 4 (Results and Discussion) evaluates the empirical performance of the adapted framework on human cell line and mouse tissue datasets. It presents clustering quality assessments using the Global Differentiation Index, identifies top candidate transcripts demonstrating differential usage, and explores biological interpretations alongside current methodological limitations.

Chapter 5 (Conclusions) summarizes the primary contributions of the work, evaluates the overall success of extending COTAN to transcript-level data, and outlines future computational directions, including long-read data integration and trajectory-aware analysis.
